\section{Reduction to Cyclic Torsors on \texorpdfstring{$\bP^1$}{P1}}
\label{section:reductions_before_cyclic_theories}

Before introducing explicit equations, we explain the strategy of the proof
and isolate two reductions that simplify the situation. The theorem concerns
a torsor under an arbitrary finite abelian group on an arbitrary smooth curve,
whereas the calculation we ultimately wish to perform is local and cyclic. We
show that both the structure group and the geometry may be simplified without
changing the ramification problem relevant to the exceptional divisor.

The first reduction is algebraic: using the decomposition of a finite abelian
group into cyclic factors, together with the compatibility of the
symmetric-power construction and ramification bounds with products and
quotients, we reduce the theorem to cyclic groups of order prime to $p$ and
cyclic groups of order $p^r$. This explains why the three cyclic theories
recalled in the next section suffice. The second reduction is geometric and
formal--local: the ramification of the induced torsor at the generic point of
the exceptional divisor depends only on the restriction of the original
torsor to the punctured formal neighbourhood of the point supporting the
modulus. It therefore tells us exactly what the cyclic theories must
accomplish: realize the resulting local cyclic torsor by an explicit equation
on $\bG_m \subset \bP^1_k$.

Together, these reductions transform the original problem into three parallel
explicit calculations, corresponding to Kummer, Artin--Schreier, and
Artin--Schreier--Witt theory. The following section recalls precisely the
conclusions from these theories that provide the required realizations.

\subsection{Reduction from Finite Abelian Groups to Cyclic Groups}

Fix a single-point submodulus $\ndls$ and write $d=\deg(\ndls)$. The first
reduction concerns the structure group and is independent of the geometry of
the curve.

\begin{prop}[Reduction to cyclic groups]
\label{prop:reduction_to_cyclic_groups}
Suppose that the conclusion of
\Cref{theorem:torsors_after_blowup_is_tame} is known whenever the structure
group is cyclic of order prime to $p$, and whenever the structure group is a
cyclic $p$-group. Then it holds for every finite abelian group $G$.
\end{prop}

\begin{proof}
By the structure theorem for finite abelian groups, choose a decomposition
\[
G \cong
\prod_{i=1}^s \Z/e_i\Z \times \prod_{j=1}^t \Z/p^{r_j}\Z,
\qquad (e_i,p)=1.
\]
Let $\rho:\pi_1^{\text{\'et}}(U,\bar{x})\to G$ be the character
corresponding to $\cP$. Composing $\rho$ with the projections to the displayed
cyclic factors gives torsors $\cP_i$ under $\Z/e_i\Z$ and torsors $\cQ_j$
under $\Z/p^{r_j}\Z$. Since the character of each factor is a quotient of
$\rho$, the ramification bound on $\cP$ passes to every $\cP_i$ and $\cQ_j$.

The original torsor is recovered as their fiber product over $U$:
\[
\cP \cong
\cP_1 \times_U \cdots \times_U \cP_s
\times_U \cQ_1 \times_U \cdots \times_U \cQ_t.
\]
The construction of the symmetric-power torsor commutes with this product.
Indeed, after pullback to $U^d$, both sides are obtained by taking the
contracted product of the $d$ pullbacks along multiplication in the structure
group. Hence
\[
\cP^{[d]} \cong
\cP_1^{[d]} \times_{U^{(d)}} \cdots \times_{U^{(d)}} \cP_s^{[d]}
\times_{U^{(d)}} \cQ_1^{[d]} \times_{U^{(d)}} \cdots
\times_{U^{(d)}} \cQ_t^{[d]}.
\]
By the assumed cyclic cases, every factor on the right is tamely ramified at
$\eta_\ndls$. The fiber product is therefore tamely ramified there by
\Cref{lemma:ramification_of_fiber_product}. This proves the assertion for
$G$.
\end{proof}

Thus it remains to understand cyclic groups. Their prime-to-$p$ part is
described by Kummer theory, while cyclic groups of order $p$ and $p^r$ are
described by Artin--Schreier and Artin--Schreier--Witt theory, respectively.

\subsection{Formal--Local Reduction from \texorpdfstring{$C$}{C} to \texorpdfstring{$\bP^1$}{P1}}

We argue that it is enough to prove \Cref{theorem:torsors_after_blowup_is_tame} for $C=\bP^1_k$, $\ndls = d \cdot 0$, and $\cP \to \bG_m = \bP^1_k \setminus \{0, \infty\}$ a $G$-torsor given by an explicit global equation.
Informally, this is because \emph{the ramification of $\cP^{[d]}$ at $\eta_\ndls$ is a formal--local invariant of $\cP$ at $P$}:
it depends only on the completed local ring $\widehat{\Oo_{C,P}} \cong k[[x]]$ together with the restriction of $\cP$
to the punctured formal disc $\Spec(k((x)))$, since all of the constructions involved --- self products, quotients by finite
groups, and blowups --- commute with the flat base change to the completion; for blowups, this is
\Cref{theorem:blowup_flat_base_change}.
We emphasize that this is not a general principle that can simply be quoted: it is a statement about the particular
constructions at hand, and it is proved by checking each construction separately. This is the content of
\Cref{lemma:formal_local_invariance} and \Cref{cor:reduction_to_P1} below.
The second ingredient of the reduction is \Cref{lemma:realization_over_Gm}: every local torsor class appearing in
\Cref{theorem:torsors_after_blowup_is_tame} is realized by an explicit torsor over $\bG_m \subset \bP^1_k$.

Throughout this section we assume that $k$ is algebraically closed; in particular every closed point of $C$ is $k$-rational, and $\ndls = dP$ with $P \in C(k)$ and $d = \deg \ndls$.

\begin{rem*}
    The assumption that $k$ is algebraically closed is harmless in our situation, for the following reasons.
    \begin{enumerate}
        \item \textbf{(Base change to $\bar{k}$.)} Assume $k$ is perfect, so that $\bar{k}=k^{sep}$.
        The formation of $C^{(d)}$ (a quotient by a finite group), of the blowup $X_\ndls$ along $Z_\ndls$, and of
        $\cP^{[d]}$ (\Cref{prop:symmetric_power_semidirect_quotient}) commutes with the flat base change $\Spec(\bar{k}) \to \Spec(k)$.
        Since $E_\ndls \cong \bP^{\deg \ndls - 1}_k$ is geometrically irreducible, $(E_\ndls)_{\bar{k}}$ is irreducible
        and its generic point $\bar{\eta}$ lies over $\eta_\ndls$; moreover
        $\Oo_{X_\ndls, \eta_\ndls} \to \Oo_{(X_\ndls)_{\bar{k}}, \bar{\eta}}$ is a weakly unramified extension of discrete
        valuation rings with separable residue field extension. Therefore, by
        \Cref{lemma:weakly_unramified_base_change},
        $(\cP_{\bar{k}})^{[d]}=(\cP^{[d]})_{\bar{k}}$ is unramified
        (resp.\ tamely ramified) at $\bar{\eta}$ if and only if
        $\cP^{[d]}$ is unramified (resp.\ tamely ramified) at $\eta_\ndls$.
        Thus both properties may be checked after base change to $\bar{k}$.
        \item \textbf{(Rationality of $P$ --- a caveat.)} If $P$ is a closed point with $\kappa(P) \neq k$, then after base
        change to $\bar{k}$ the point $Z_{\ndls}$ corresponds to the divisor $\sum_{\bar{P} \mapsto P} n \bar{P}$ on
        $C_{\bar{k}}$, whose support consists of $[\kappa(P):k]$ distinct points, whereas the computation below treats a
        modulus supported at a single rational point.
        {\color{red}
        We therefore prove \Cref{theorem:torsors_after_blowup_is_tame} only
        under the standing assumption that $k$ is algebraically closed. This
        is the version used later: at the beginning of the proof of
        \Cref{theorem:SymmetricPowerOfSheafIsTamelyRamified}, all data are
        first base changed to $\bar{k}$. The modulus then decomposes into its
        rational-point-supported summands; the present theorem is applied to
        each summand, and \Cref{lemma:moduli_reduction} combines them before
        tameness is descended back to $k$.
        }
    \end{enumerate}
\end{rem*}

We now formulate the formal--local invariance precisely. Fix $P \in C(k)$ lying in the support of $\mdls$, set $\ndls = dP$,
and fix a uniformizer $x$ of $\Oo_{C,P}$; it induces an isomorphism $\widehat{\Oo_{C,P}} \cong k[[x]]$.
The scheme $\Spec(k[[x]])$ has two points --- the closed point, mapping to $P$, and the generic point, mapping to the generic
point of $C$ --- so the punctured formal disc $\Spec(k((x)))$ maps to $U$, and we denote by
$$\cP_x := \cP \times_U \Spec(k((x)))$$
the corresponding restriction of $\cP$; it is a $G$-torsor over $\Spec(k((x)))$.
Let $e_1, \dots, e_d \in k[[x_1, \dots, x_d]]$ be the elementary symmetric polynomials in $x_1, \dots, x_d$, and set
$$ V = \Spec\left(k[[e_1, \dots, e_d]][e_d^{-1}]\right), \qquad W = \Spec\left(k[[x_1, \dots, x_d]][(x_1 \cdots x_d)^{-1}]\right).$$
The inclusion $k[[e_1, \dots, e_d]] \subset k[[x_1, \dots, x_d]]$ induces a finite flat morphism $\hat{r}: W \to V$
(note that inverting $e_d = x_1 \cdots x_d$ inverts each $x_i$), and for each $i$ we let
$$ q_i : W \to \Spec(k((x))), \qquad x \mapsto x_i.$$
The group $S_d$ acts on $W$ over $V$ by permuting $x_1, \dots, x_d$, and $q_{\sigma(i)} = q_i \circ \sigma$.
Recall from \Cref{prop:symmetric_power_semidirect_quotient} that
\[
K_G:=\ker\bigl(G^d\xrightarrow{\,m\,}G\bigr),
\qquad
\Xi_G:=K_G\rtimes S_d,
\]
where $m(g_1,\dots,g_d)=g_1\cdots g_d$ and $S_d$ acts on $K_G$ by
permuting coordinates. Thus quotienting by $\Xi_G$ means first contracting
the $G^d$-torsor along multiplication, by quotienting by $K_G$, and then
descending from the ordered product to the symmetric product by quotienting
by $S_d$.

\begin{lemma}[Formal--local invariance]\label{lemma:formal_local_invariance}
    In the above notation:
    \begin{enumerate}
        \item The choice of $x$ induces an isomorphism $\widehat{\Oo_{C^{(d)}, Z_\ndls}} \cong k[[e_1, \dots, e_d]]$ carrying
        the maximal ideal to $(e_1, \dots, e_d)$, such that the induced morphism $c: \Spec(k[[e_1, \dots, e_d]]) \to C^{(d)}$
        is flat and satisfies $c^{-1}(U^{(d)}) = V$; we write $c_V : V \to U^{(d)}$ for the restriction.
        \item There is an induced isomorphism $\widehat{\Oo_{X_\ndls, \eta_\ndls}} \cong \kappa[[e_d]]$, where
        $\kappa = k(\frac{e_1}{e_d}, \dots, \frac{e_{d-1}}{e_d})$. In particular
        $\widehat{K}_{\eta_\ndls} := \Frac(\widehat{\Oo_{X_\ndls, \eta_\ndls}}) \cong \kappa((e_d))$, and the canonical
        morphism $\Spec(\widehat{K}_{\eta_\ndls}) \to U^{(d)}$ factors as
        $\Spec(\kappa((e_d))) \to V \xrightarrow{c_V} U^{(d)}$, where the first map does not depend on $(C, \cP)$.
        \item There is a canonical isomorphism of $G$-torsors over $V$
        $$ \cP^{[d]} \times_{U^{(d)}} V \;\cong\; \Xi_G \backslash \left( q_1^{-1}\cP_x \times_W \cdots \times_W q_d^{-1}\cP_x \right).$$
    \end{enumerate}
    Consequently, the isomorphism class of the $G$-torsor
    $\cP^{[d]}|_{\Spec(\widehat{K}_{\eta_\ndls})}$ depends only on $d$ and on
    the isomorphism class of the $G$-torsor $\cP_x$ over $k((x))$, and not
    otherwise on the pair $(C,\cP)$. In particular, whether $\cP^{[d]}$ is
    tamely ramified or unramified at $\eta_\ndls$ depends only on
    $(\cP_x,d)$.
\end{lemma}

\begin{proof}
    \textbf{(1)} The chosen parameter gives
    $\widehat{\Oo_{C^d,(P,\dots,P)}}\cong k[[x_1,\dots,x_d]]$.
    We use the standard facts that completion is compatible with finite
    morphisms and that the completed local ring of a finite quotient is the
    invariant ring of the completed local ring.  Since $(P,\dots,P)$ is the
    unique point of $C^d$ above $dP$, this gives
    \[
        \widehat{\Oo_{C^{(d)},dP}}
        \cong k[[x_1,\dots,x_d]]^{S_d}
        =k[[e_1,\dots,e_d]],
    \]
    where the last equality is the formal-power-series version of the
    fundamental theorem of symmetric functions.  The maximal ideal is thus
    $(e_1,\dots,e_d)$, and the resulting map $c$ is flat by flatness of
    completion.

    It remains only to identify the boundary.  Near $dP$, the only component
    of $C^{(d)}\setminus U^{(d)}$ that occurs is the locus of divisors
    containing $P$.  Its inverse image in the ordered formal neighbourhood is
    \[
        \bigcup_i V(x_i)=V(x_1\cdots x_d)=V(e_d).
    \]
    Since the ordered-to-symmetric map is finite and surjective, the boundary
    in $\Spec k[[e_1,\dots,e_d]]$ is also $V(e_d)$.  Hence
    $c^{-1}(U^{(d)})=V$ (and similarly the corresponding ordered inverse image
    is $W$).

    \textbf{(2)} We use that blowups commute with flat base change.  Thus the
    pullback of $X_\ndls$ along $c$ is the blowup of
    $\Spec k[[e_1,\dots,e_d]]$ at its closed point.  Completion preserves the
    exceptional divisor and its residue field, so it is enough to compute the
    completed local ring at the generic point of the exceptional divisor in
    this standard blowup.

    On the chart where $e_d$ is the exceptional parameter, put
    $t_i=e_i/e_d$ for $i<d$.  The exceptional divisor is given by $e_d=0$,
    and its generic point has residue field
    \[
        \kappa=k(t_1,\dots,t_{d-1})
        =k\left(\frac{e_1}{e_d},\dots,\frac{e_{d-1}}{e_d}\right).
    \]
    The usual local calculation for the blowup of a regular local scheme at
    its closed point therefore gives
    $\widehat{\Oo_{X_\ndls,\eta_\ndls}}\cong\kappa[[e_d]]$ and
    $\widehat K_{\eta_\ndls}\cong\kappa((e_d))$.  Since $e_d$ is invertible
    in the latter field, its map to the completed symmetric power factors
    through $V$, giving the asserted factorization through $c_V$.

    \textbf{(3)} Let $h : \cP^{\times_k d} \to U^{(d)}$ denote the canonical morphism; it is finite, being the composition
    of the finite étale morphism $\cP^{\times_k d} \to U^d$ with the finite morphism $U^d \to U^{(d)}$. By
    \Cref{prop:symmetric_power_semidirect_quotient},
    $\cP^{[d]} = \Xi_G \backslash \cP^{\times_k d} = \underline{\Spec}_{U^{(d)}}\left((h_* \Oo_{\cP^{\times_k d}})^{\Xi_G}\right)$.
    Pushforward along the affine morphism $h$ commutes with the flat base change $c_V : V \to U^{(d)}$, and the formation of
    $\Xi_G$-invariants is a kernel, hence also commutes with flat base change; therefore
    $$ \cP^{[d]} \times_{U^{(d)}} V \cong \Xi_G \backslash \left(\cP^{\times_k d} \times_{U^{(d)}} V\right).$$
    Completing the finite map $C^d\to C^{(d)}$ at the unique point
    $(P,\dots,P)$ above $dP$ gives the Cartesian square defined by
    $k[[e_1,\dots,e_d]]\subset k[[x_1,\dots,x_d]]$.  Restricting this square
    to $e_d=x_1\cdots x_d\ne0$ gives directly
    $U^d\times_{U^{(d)}}V\cong W$, and hence
    $\cP^{\times_k d}\times_{U^{(d)}}V
      \cong\cP^{\times_k d}\times_{U^d}W$.

    The $i$-th component of the resulting map $W\to U^d$ is
    $W\xrightarrow{q_i}\Spec(k((x)))\to U$: indeed, it is induced in
    coordinates by $x\mapsto x_i$, and $x_i$ is invertible on $W$.
    Since
    $\cP^{\times_k d}=p_1^{-1}\cP\times_{U^d}\cdots\times_{U^d}p_d^{-1}\cP$,
    compatibility of fiber products with base change now gives
    $$ \cP^{\times_k d} \times_{U^d} W \cong q_1^{-1}\cP_x \times_W \cdots \times_W q_d^{-1}\cP_x, $$
    which is equivariant for the actions of $G^d$ (in particular of $\KG$) and of $S_d$ --- the latter acting on $W$ by
    permuting the $x_i$, compatibly with $q_{\sigma(i)} = q_i \circ \sigma$. Passing to $\Xi_G$-quotients yields the displayed
    isomorphism.

    For the final statement: the schemes $V$ and $W$, the maps $q_i$ and $\hat{r}$, the $\Xi_G$-action, and the point
    $\Spec(\kappa((e_d))) \to V$ of (2) are all constructed from $d$ alone; by (3), the torsor
    $\cP^{[d]}|_{\Spec(\widehat{K}_{\eta_\ndls})}$ is the pullback along $\Spec(\kappa((e_d))) \to V$ of a torsor constructed
    from $\cP_x$ and these data. Since tame ramification and unramifiedness at
    $\eta_\ndls$ are, by definition, determined by this restricted torsor
    together with the discrete valuation ring $\kappa[[e_d]]$
    (\Cref{definition:tamely_ramified_in_codim_1}), both properties depend only
    on $(\cP_x,d)$.
\end{proof}

\begin{cor}[Reduction to $\bP^1$]\label{cor:reduction_to_P1}
    Let $C, C'$ be smooth projective geometrically connected curves over the algebraically closed field $k$, let
    $\cP \to U = C \setminus \mdls$ and $\cP' \to U' = C' \setminus \mdls'$ be $G$-torsors, let
    $P \in \mathrm{Supp}(\mdls)(k)$ and $P' \in \mathrm{Supp}(\mdls')(k)$, and let $x, x'$ be uniformizers at $P, P'$
    respectively. Assume that $\cP_x \cong \cP'_{x'}$ as $G$-torsors, compatibly with the isomorphism
    $k((x)) \cong k((x'))$, $x \mapsto x'$.
    Then, for every $d\geq1$, $\cP^{[d]}$ is unramified
    (resp.\ tamely ramified) at $\eta_{dP}$ if and only if
    $\cP'^{[d]}$ is unramified (resp.\ tamely ramified) at $\eta_{dP'}$.
\end{cor}
\begin{proof}
    Immediate from \Cref{lemma:formal_local_invariance}: both $\cP^{[d]}|_{\Spec(\widehat{K}_{\eta_{dP}})}$ and
    $\cP'^{[d]}|_{\Spec(\widehat{K}_{\eta_{dP'}})}$ are identified, compatibly with the discrete valuation ring
    $\kappa[[e_d]]$, with the pullback along $\Spec(\kappa((e_d))) \to V$ of
    $\Xi_G \backslash (q_1^{-1}\cP_x \times_W \cdots \times_W q_d^{-1}\cP_x)$.
\end{proof}

The corollary reduces the geometry to $\bP^1$ once the local cyclic torsor is
realized by a global equation on $\bG_m$. The next section introduces the three
cyclic theories that provide precisely these equations and control their
ramification.
